\section{Related Work}

\subsection{Autoregressive Language Models}
Autoregressive language models \citep{radford2019language,touvron2023llama,yang2025qwen3,liu2024deepseek} factorize the discrete text distribution by the chain rule and are trained with token-level maximum likelihood, making them the most widely adopted paradigm for text modeling. Their limitations are that generation is constrained by a fixed left-to-right order, inference is inherently sequential, and they are less suitable for non-monotonic generation tasks such as infilling, local editing, and global reorganization. In contrast, \method first models a global semantic prior in a continuous latent space and then performs conditional decoding, thereby alleviating token-level ordering bias and improving generation efficiency with a block-causal DiT.

\subsection{Discrete Diffusion Language Models}
Discrete diffusion language models mainly fall into two categories. The first category is based on discrete transition kernels \citep{austin2021structured,campbell2022continuous,sun2022score}, which define forward perturbation and reverse recovery in discrete token space and achieve non-autoregressive generation through multi-step denoising; however, sampling is usually slow and these methods cannot easily exploit the smooth semantic structure of continuous spaces. The second category is based on masking or absorbing states \citep{sahoo2024simple,nie2025large,zhu2025llada,you2025llada,zhu2025lladamoe,shi2024simplified,zheng2025continuously,zheng2024masked,sahoo2025diffusion,nie2024scaling}, which construct training objectives by progressively mapping tokens to masks or absorbing states and then recovering the original text; however, information loss in intermediate states limits global semantic planning and fine-grained control. In contrast, \method moves the diffusion process to a continuous latent space, where compressible latent variables carry global semantics, thus combining the manipulability of continuous spaces with hierarchical semantic modeling.

\subsection{Continuous Diffusion Language Models}
Continuous diffusion language models can be broadly divided into three categories. The first category consists of high-dimensional vocabulary-aligned continuous methods \citep{han2023ssd,richemond2022categorical,mahabadi2024tess,jo2025continuous}, which perform continuous diffusion or flow modeling directly on one-hot vectors, logit simplexes, or probability simplexes to preserve alignment with discrete vocabularies; however, their representation dimension scales with vocabulary size, which limits scalability. The second category consists of token-embedding-based continuous methods \citep{li2022diffusion,strudel2022self,gao2024empowering,gong2022diffuseq,gulrajani2023likelihood,chen2023cheaper,lin2023text,dieleman2022continuous}, which first map text into continuous embedding spaces and then apply diffusion or flow modeling to improve generation flexibility; however, their generation process remains essentially the recovery of noisy target representations, lacking an explicit hierarchical latent-variable interpretation and a unified marginal-likelihood view of text distributions. The third category consists of latent-space continuous methods \citep{meshchaninov2025cosmos,kang2025ladir,lovelace2023latent,zhang2023planner}, which compress text into latent spaces with autoencoders or VAEs and then perform diffusion modeling. These methods typically rely on latent-space design and autoregressive decoders, and usually treat the latent space as a fixed representation rather than modeling it under a hierarchical latent-variable framework. In contrast, \method explicitly separates global semantics from local realization through hierarchical latent-variable modeling, and learns a semantic prior in a dynamic continuous latent space, thereby better balancing modeling flexibility, inference efficiency, and theoretical interpretability.